%-------------------------
% Anvi Wadhwa — Resume
%-------------------------

\documentclass[10pt,letterpaper]{article}

\usepackage[
  top=0.35in,
  bottom=0.35in,
  left=0.45in,
  right=0.45in
]{geometry}

\usepackage{enumitem}
\usepackage{titlesec}
\usepackage{hyperref}
\usepackage[dvipsnames,usenames]{xcolor}
\usepackage[T1]{fontenc}
\usepackage{tabularx}
\usepackage{array}

\hypersetup{
  colorlinks=true,
  urlcolor=blue!60!black,
  linkcolor=blue!60!black
}

\pagestyle{empty}
\setlength{\parindent}{0pt}
\setlength{\parskip}{0pt}

% Section styling
\definecolor{accent}{RGB}{26,60,107}

\titleformat{\section}
  {\vspace{2pt}\bfseries\small\color{accent}\uppercase}
  {}{0em}{}
  [\vspace{-2pt}\color{accent}\rule{\linewidth}{0.6pt}\vspace{2pt}]

\titlespacing{\section}{0pt}{4pt}{3pt}

% Bullet styling
\setlist[itemize]{
  leftmargin=1.2em,
  itemsep=0.8pt,
  parsep=0pt,
  topsep=1pt,
  label=\textbullet
}

% Role header (NO ITALICS)
\newcommand{\roleheader}[3]{%
  \textbf{#1} \hfill {\small\color{gray}#2}\\
  {\small\color{gray}#3}
}

% Project header (NO ITALICS)
\newcommand{\projheader}[3]{%
  \textbf{#1} \hfill {\small\color{gray}#2}%
  \ifx&#3&\else\\{\small\color{gray}#3}\fi
}

\begin{document}

%──────────────────────────── HEADER ────────────────────────────

\begin{center}
{\LARGE\bfseries Anvi Wadhwa}\\[3pt]

{\small
+91\,93190\,86177 \;$\cdot$\;
\href{mailto:email@example.com}{anvi.wadhwa2024@vitstudent.ac.in}
\;$\cdot$\;
\href{https://github.com/anvisgit?tab=repositories}{github}
\;$\cdot$\;
\href{https://linkedin.com/in/anviwadhwa7}{linkedin}
\;$\cdot$\;
\href{https://medium.com/@anviwadhwa7}{medium.com/@anviwadhwa7}
\;$\cdot$\;
\href{https://anvisgit.github.io/anvisgit/}{Portfolio Website}
}
\end{center}

\vspace{-4pt}

%────────────────────── RESEARCH INTERESTS ─────────────────────

\section{Research Interests}

{\small
I like to apply machine learning to physical systems, focusing on modeling interplanetary magnetic field fluctuations from muon-based cosmic-ray data. I have worked with neural networks and graph neural networks for high-dimensional structure learning, Bayesian methods for probabilistic inference, and reinforcement learning, with emphasis on optimization dynamics, convergence behavior, and uncertainty quantification.
}

%──────────────────────── EDUCATION ────────────────────────────

\section{Education}

\roleheader
{Vellore Institute of Technology (VIT Chennai)}
{Aug 2024 -- May 2028}
{B.Tech, Computer Science \quad$\cdot$\quad CGPA: 8.1/10 \quad$\cdot$\quad Chennai, India}

\vspace{1pt}

\begin{itemize}
\item R\&D Lead, Business Innovation Community; directed applied AI/ML projects for a 12-member research team
\item Active member of Microsoft Innovation Community and Hackclub, contributing to AI/ML and blockchain initiatives
\end{itemize}

%──────────────────────── EXPERIENCE ───────────────────────────

\section{Experience}

\roleheader
{Computer Vision Research Intern --- IIT (BHU) Varanasi}
{May 2025 -- Present}
{AI-Powered Rehabilitation \& Mobility System}

\vspace{1pt}

\begin{itemize}
\item Developed a real-time physiotherapy assessment system using MediaPipe and OpenCV, tracking lower-limb joint kinematics at \textbf{28 FPS} on CPU-only hardware
\item Engineered a joint-angle estimation pipeline achieving \textbf{$\pm$2° accuracy} across 6 DOF, reducing reliance on manual goniometric measurements

\end{itemize}

%────────────────────── -------------- PROJECTS ──────────────────────

\section{Projects}

\projheader
{Comparative Analysis of PPO and GRPO for Solving Wordle| Building }
{}

\begin{itemize}
\item Implemented PPO and GRPO from scratch in a custom Gymnasium Wordle environment and benchmarked training dynamics under identical hyperparameter settings
\end{itemize}

\vspace{2pt}

\projheader
{Systematic Intraday Trading via GARCH Volatility Forecasting}
{}

\begin{itemize}
\item Developed a systematic intraday trading framework combining rolling GARCH(1,3) volatility forecasts with RSI-Bollinger mean-reversion signals on high-frequency equity data
\item Evaluated strategy robustness using walk-forward validation and out-of-sample testing, achieving \textbf{+5.0\% test return}, \textbf{0.65 Sharpe ratio}, and \textbf{-6.6\% maximum drawdown}
\end{itemize}

\vspace{2pt}

\projheader
{Speech Emotion Recognition-- MultiBranch 1D CNN}
{}

\begin{itemize}
\item Designed a multi-branch 1D-CNN architecture combining MFCC, Mel spectrogram, and Chroma representations through late-fusion aggregation
\item Achieved \textbf{75.7\% accuracy} and \textbf{75.5\% macro F1} on the RAVDESS emotion-recognition benchmark
\item Improved performance by \textbf{7.5 percentage points} over an MFCC-only baseline, demonstrating the value of spectral feature diversity
\end{itemize}

%────────────────────── TECHNICAL SKILLS ───────────────────────

\section{Technical Skills}

\begin{tabular}{@{}ll}
\small\textbf{Programming} &
\small Python, Java \\[2pt]

\small\textbf{Machine Learning} &
\small Numpy, Pandas, PyTorch, TensorFlow, Scikit-learn, ARCH/GARCH Models\\[2pt]

\small\textbf{Reinforcement Learning} &
\small PPO, GRPO, Stable-Baselines3, Gymnasium, Policy Gradient Methods \\[2pt]

\small\textbf{Computer Vision \& DSP} &
\small OpenCV, MediaPipe, MoveNet(TF) \\[2pt]

\small\textbf{Tools} &
\small Git, Linux
\end{tabular}

%────────────────────── CERTIFICATIONS ─────────────────────────

\section{Certifications}

Machine Learning Specialization --- Coursera / Andrew Ng |  Quantitative Finance --- IIT Kharagpur


\end{document}

